\section{Testing und Qualitätssicherung}
\label{sec:testing}

\subsection{Teststrategie}
\label{sec:teststrategie}

Für jede Schwachstelle galt dieselbe Arbeitsregel: erst beweisen, dass der
Fehler wirklich existiert. Dann erst beheben. Das klingt offensichtlich, ist in
der Praxis aber entscheidend: Ein Test, der vor dem Patch scheitert und danach
besteht, ist der einzige belastbare Nachweis, dass die Lücke geschlossen ist.
Das Vorgehen gliedert sich in fünf Schritte:

\begin{enumerate}
      \item \textbf{Einlesen und verstehen.}
            Das gemeldete Issue lesen, den betroffenen Code durchgehen --
            \texttt{base.py}, Modell-Definitionen, Router, bis der
            Angriffspfad klar ist.
      \item \textbf{Fehler reproduzieren.}
            Einen minimalen HTTP-Request zusammenstellen, der die Schwachstelle
            auslöst (etwa GET auf eine fremde Ressource mit eigenem Token).
            Vor dem Patch antwortet der Server mit 200, obwohl er 404
            zurückgeben müsste.
      \item \textbf{Failing Test schreiben.}
            Den reproduzierten Angriff als automatisierten Test kodieren.
            Auf dem ungepatchten Stand muss der Test scheitern (\emph{Red}).
      \item \textbf{Patch implementieren.}
            Die kleinste Code-\u00c4nderung vornehmen, die den Test bestehen lässt.
      \item \textbf{Tests ausführen.}
            Alle neuen \emph{und} bestehenden Tests laufen lassen. Neue Tests
            grün; kein bestehender darf neu scheitern (\emph{Green}).
\end{enumerate}

\noindent Tests werden in drei Kategorien geführt:

\begin{itemize}
      \item \textbf{Sicherheitstests} (\texttt{tests/test\_idor.py}): Eine einzelne
            Datei, die Proof-of-Concept und Regressionstests vereint.
            Deckt beide Angriffsklassen über alle relevanten HTTP-Verben ab
            (GET, PATCH, DELETE, LIST).
      \item \textbf{Bestehende Suite} (Smoke- und Integrationstests):
            Stellen sicher, dass der Patch keine reguläre Funktionalität bricht.
\end{itemize}

\subsection{Neue Sicherheitstests}
\label{sec:neue-sicherheitstests}

Alle Sicherheitstests befinden sich in \texttt{tests/test\_idor.py}.
Jeder Test registriert zwei Mandanten über die öffentliche
\texttt{/auth/register}-API und erstellt Ressourcen via \texttt{POST /vehicles}
bzw. \texttt{POST /employees}, ohne direkt auf die Datenbank zuzugreifen.
Die Datei deckt zwei Angriffsklassen über mehrere HTTP-Verben ab:

\paragraph{Klasse 1: IDOR (Issue \#5).}
Zwei Organisationen (Org A = Opfer, Org B = Angreifer) werden angelegt.
Org B versucht, Ressourcen von Org A zu lesen, zu ändern und zu löschen:

\begin{center}
      \begin{tabular}{|l|l|}
            \hline
            \textbf{Testfall}                 & \textbf{Erwartetes Ergebnis}            \\
            \hline
            \texttt{test\_get\_blocked}       & GET fremdes Objekt → 404                \\
            \texttt{test\_patch\_blocked}     & PATCH fremdes Objekt → 404              \\
            \texttt{test\_delete\_no\_effect} & DELETE: Objekt bei Org A unverändert    \\
            \texttt{test\_list\_isolated}     & LIST enthält keine Fremd-Daten          \\
            \texttt{test\_no\_org\_leak}      & Antwortbody enthält keine fremde Org-ID \\
            \hline
      \end{tabular}
\end{center}

\vspace{0.5em}
\noindent HTTP-DELETE ist per Standard idempotent; der Statuscode ist 204,
unabhängig davon, ob ein Objekt gefunden wurde. Der bestehende Smoke-Test \texttt{test\_delete\_vehicle\_nonexistent\_returns\_204}
bestätigt diese Eigenschaft. Das sicherheitskritische Merkmal ist daher nicht
der Statuscode, sondern dass Org A das Objekt anschliessend noch lesen kann.

\paragraph{Klasse 2: Soft-Delete-Bypass (Issue \#14).}
Eine Ressource wird soft-gelöscht. Alle nachfolgenden Zugriffe auf sie müssen
\texttt{404} zurückgeben:

\begin{center}
      \begin{tabular}{|l|l|}
            \hline
            \textbf{Testfall}             & \textbf{Erwartetes Ergebnis}     \\
            \hline
            \texttt{test\_get\_gone}      & GET soft-deleted → 404           \\
            \texttt{test\_patch\_gone}    & PATCH soft-deleted → 404         \\
            \texttt{test\_list\_hidden}   & LIST filtert soft-deleted heraus \\
            \texttt{test\_double\_delete} & Zweites DELETE → 204, GET → 404  \\
            \hline
      \end{tabular}
\end{center}

\paragraph{Klasse 3: Modulübergreifende Verifizierung (\texttt{TestIDOREmployee}).}
Der \texttt{BaseService}-Patch gilt für alle erbenden Services.
Um dies empirisch zu belegen, wiederholt \texttt{TestIDOREmployee} dieselben
Angriffsszenarien mit \emph{Mitarbeitenden} als Zielressource:

\begin{center}
      \begin{tabular}{|l|l|}
            \hline
            \textbf{Testfall}                 & \textbf{Erwartetes Ergebnis}           \\
            \hline
            \texttt{test\_get\_blocked}       & GET fremdes Objekt → 404               \\
            \texttt{test\_patch\_blocked}     & PATCH fremdes Objekt → 404             \\
            \texttt{test\_list\_isolated}     & LIST enthält keine Fremd-Einträge      \\
            \texttt{test\_soft\_delete\_gone} & GET/LIST nach DELETE → 404 bzw. absent \\
            \hline
      \end{tabular}
\end{center}

\subsection{Ergänzende Schwachstellenanalyse (Phase 2 \& 3)}
\label{sec:zusatz-schwachstellen}

\subsubsection{Warum erst nach Phase 1?}

Der initiale Patch in \texttt{base.py} behebt IDOR und Soft-Delete-Bypass
vollständig, aber \emph{nur für Code, der durch \texttt{BaseService} läuft}.
Sobald die Phase-1-Tests grün waren, stellte sich die Folgefrage:

\begin{quote}
      \emph{Gibt es Endpunkte, die \texttt{BaseService} umgehen und daher vom zentralen
            Patch unberührt bleiben?}
\end{quote}

Die Antwort erforderte eine gezielte Durchsicht aller Router-Dateien, ein Schritt,
der sinnvollerweise erst \emph{nach} dem Basis-Fix stattfindet,
weil vorher der Patch-Scope noch nicht stabil war.

\subsubsection{Systematische Vorgehensweise}

Für jede Router-Datei wurde die folgende Checkliste angewendet:

\begin{enumerate}
      \item Welche Service-Methode ruft der Endpunkt auf?
      \item Führt diese Methode zu \texttt{BaseService.get()} oder
            \texttt{get\_multi()}?
            Falls \textbf{nein}: Enthält die Methode einen eigenen
            \texttt{organization\_id}-Filter?
      \item Wird nach einem \texttt{get\_or\_404}-Lookup die
            \texttt{user\_id} oder \texttt{organization\_id} des Ergebnisses
            gegen den eingeloggten Benutzer geprüft?
\end{enumerate}

\noindent Diese Checkliste wurde in zwei Durchgängen angewendet:
\textbf{Phase 2} nach dem Basis-Patch (Dokumente, SAT, Notifications),
\textbf{Phase 3} nach der vollständigen Erweiterung auf alle verbleibenden Module.

\subsubsection{Phase-2-Findings (3 Lückenklassen)}

\begin{enumerate}
      \item \textbf{Notifications: fehlende Eigentümerprüfung (GET / DELETE / mark-as-read).}\\
            \texttt{get\_or\_404()} liefert die Notification zurück, aber ein
            Vergleich \texttt{notification.user\_id == current\_user.id} fehlte
            in \emph{allen drei} Endpunkten.\\
            Auswirkung: Jeder authentifizierte Benutzer kann fremde Notifications
            lesen, als gelesen markieren und löschen.\\
            \emph{Fix:} Ownership-Check \texttt{notification.user\_id != current\_user.id}
            → HTTP 404 in GET, POST /read und DELETE.

      \item \textbf{SAT-Assistances: fehlender Mandanten-Filter (GET).}\\
            \texttt{AssistanceService.get\_detail()} baut eigene SQLAlchemy-Filter
            und erbt nicht von \texttt{BaseService.get()}.
            Der \texttt{organization\_id}-Filter fehlte völlig.\\
            \emph{Fix:} Parameter \texttt{organization\_id} in
            \texttt{get\_detail()} aufgenommen; Router reicht
            \texttt{current\_user.organization\_id} weiter.

      \item \textbf{List-Endpunkte (Dokumente, Treibstoff, Verbrauchsmaterial):
                  fehlender Mandanten-Filter (LIST).}\\
            Fünf FK-List-Methoden filterten ausschliesslich nach der
            Fremdschlüssel-ID ohne \texttt{organization\_id}.\\
            \emph{Fix:} Optionaler Parameter \texttt{organization\_id} ergänzt;
            alle fünf Router-Aufrufe befüllen ihn.
\end{enumerate}

\subsubsection{Phase-3-Findings (3 HIGH + 14 MEDIUM)}

Der zweite Durchgang durch alle verbleibenden Modulordner
(\texttt{fleet/}, \texttt{tools/}, \texttt{sat/}, \texttt{shared/})
fand drei Klassen weiterer Schwachstellen:

\paragraph{HIGH: Direkte Cross-Tenant-Zugriffe.}
\begin{enumerate}
      \item \textbf{\texttt{shared/api/users.py}: GET / PATCH / DELETE
            \texttt{/{user\_id}}} (Admin-Endpunkte).\\
            \texttt{get\_or\_404} und \texttt{delete()} wurden ohne
            \texttt{organization\_id} aufgerufen.
            Ein Admin aus Org A konnte Benutzer anderer Mandanten lesen,
            ändern und löschen.\\
            \emph{Fix:} \texttt{organization\_id=current\_user.organization\_id}
            an alle drei Aufrufe übergeben.

      \item \textbf{\texttt{sat/api/attachments.py}: DELETE \texttt{/{attachment\_id}}}.\\
            \texttt{SatAttachmentService} erbt von \texttt{BaseService} und
            unterstützt \texttt{organization\_id} in \texttt{delete()},
            der Router übergab es jedoch nicht.
            Jeder SAT-Manager konnte fremde Anhänge löschen.\\
            \emph{Fix:} \texttt{organization\_id=current\_user.organization\_id}
            hinzugefügt.

      \item \textbf{\texttt{fleet/api/gps.py}: POST \texttt{/trips/\{trip\_id\}/position}}.\\
            GPS-Koordinaten konnten einem beliebigen Trip hinzugefügt werden,
            ohne zu prüfen, ob der Trip dem eigenen Mandanten gehört.
            Ein Angreifer konnte fremde Trips mit falschen Positionsdaten
            beschreiben.\\
            \emph{Fix:} \texttt{trip\_service.get\_or\_404(db, trip\_id,
                  organization\_id=current\_user.organization\_id)} vor dem Schreiben
            der Position eingefügt.
\end{enumerate}

\paragraph{MEDIUM: Enumeration via FK-Guessing (List-Methoden).}
14 Service-Methoden in sechs Dateien filterten FK-List-Abfragen ohne
\texttt{organization\_id}: GPS-Trips (by vehicle/employee, active trip, Routenpunkte,
letzte Position), Wartungsaufgaben/-zeitpläne, Treibstoffpumpen-Lieferungen,
Werkzeuge und Werkzeugzuweisungen (by employee/vehicle), SAT-Kontakte,
SAT-Maschinen und SAT-Interventionsberichte.\\
\emph{Fix (einheitliches Muster):} Pflichtparameter \texttt{organization\_id: str} (kein Defaultwert)
in jede Methode aufgenommen; damit wird ein \texttt{WHERE organization\_id = \$1}
unbedingt an die Abfrage gehängt.
Python erzwingt, dass jeder Aufrufer den Parameter explizit übergibt.
Ein versehentliches Weglassen des Mandanten-Filters führt zu einem \texttt{TypeError}
zur Laufzeit, nicht zu einem stillen Sicherheitsleck.
Alle zugehörigen Router befüllen den Parameter mit \texttt{current\_user.organization\_id}.

\subsubsection{Testnachweis (Regression)}

Da alle Phase-2- und Phase-3-Lücken durch Code-Review gefunden wurden
(kein vorheriger Testfehler), folgt der Nachweis dem gleichen
Red/Green-Prinzip wie Phase 1:
\texttt{TestIDOREmployee} wiederholt alle Angriffsszenarien mit
\emph{Mitarbeitenden} als Ressource und beweist, dass der
\texttt{BaseService}-Patch sowie die individuellen Fixes gemeinsam
wirksam sind.
Die vollständige Smoke-Test-Suite (187 Tests) wurde nach jedem
Patch-Durchgang grün ausgeführt.

\subsection{Testabdeckung des zentralen Patches}
\label{sec:patch-abdeckung}

Der Patch befindet sich ausschliesslich in \texttt{app/shared/services/base.py},
der gemeinsamen Basisklasse aller Module (Fahrzeuge, Mitarbeitende, Werkzeuge,
Treibstoff, SAT u.a.). Alle Service-Klassen erben von \texttt{BaseService} und
nutzen dieselben Methoden \texttt{get()}, \texttt{get\_or\_404()},
\texttt{get\_multi()} und \texttt{delete()} ohne eigene Überschreibungen der
Sicherheitslogik.

Daraus folgt: ein Modultest genügt, um den Patch in der Basisklasse
nachzuweisen. Die Regressionstests verwenden \emph{Fahrzeuge} als
repräsentative Ressource und prüfen alle vier betroffenen Methoden:

\begin{center}
      \begin{tabular}{|l|l|l|}
            \hline
            \textbf{Methode in base.py} & \textbf{Fix}                            & \textbf{Geprüft durch}                              \\
            \hline
            \texttt{get()}              & org\_id-Filter + soft-delete-Ausschluss & GET- und soft-delete-Tests                          \\
            \texttt{get\_or\_404()}     & Delegation an \texttt{get()}            & PATCH- und GET-Tests                                \\
            \texttt{get\_multi()}       & org\_id-Filter + soft-delete-Ausschluss & LIST-Tests (IDOR und soft-delete)                   \\
            \texttt{delete()}           & ruft \texttt{get()} mit org\_id auf     & Cross-Tenant-DELETE (Seiteneffekt: Objekt erhalten) \\
            \hline
      \end{tabular}
\end{center}

\noindent Module wie \texttt{FuelService}, die eine eigene \texttt{delete()}-Methode
besitzen (und daher separat gepatcht wurden), sind durch den bestehenden
Smoke-Test \texttt{test\_delete\_fuel\_entry\_nonexistent\_returns\_204}
abgedeckt. Dieser Test schlug vor dem Patch mit HTTP 500 fehl und besteht
nach dem Patch.

\subsection{Anpassung bestehender Tests}
\label{sec:anpassung-tests}

Nach dem Patch scheiterten zwei bestehende Tests. Beide dokumentierten
zuvor das \emph{fehlerhafte} Verhalten:

\begin{enumerate}
      \item \textbf{\texttt{test\_delete\_then\_get\_still\_returns\_soft\_deleted}}\\
            Datei: \texttt{tests/integration/test\_vehicle\_crud.py}\\
            Alter Kommentar: \emph{``soft-deleted records are still returned by GET''}\\
            Alte Assertion: \texttt{assert get\_res.status\_code == 200}\\
            Neue Assertion: \texttt{assert get\_res.status\_code == 404}\\
            Der Testname bleibt unverändert als historisches Artefakt, der Kommentar
            wurde korrigiert.

      \item \textbf{\texttt{test\_delete\_fuel\_entry\_nonexistent\_returns\_204}}\\
            Datei: \texttt{tests/smoke/test\_fuel.py}\\
            \texttt{FuelService} besitzt eine eigene \texttt{delete()}-Methode,
            die \texttt{organization\_id} nicht als Parameter kannte.
            Der API-Endpunkt übergab den Parameter, der Override ignorierte ihn; Python warf einen \texttt{TypeError}, der zum HTTP-500 wurde.\\
            Fix: Signatur um \texttt{*, organization\_id: str} ergänzt
            und an den internen \texttt{self.get()}-Aufruf weitergeleitet.
\end{enumerate}

\subsection{Ausführung der Integrationstests}
\label{sec:testausfuehrung}

Das finale Ergebnis nach allen Korrekturen:

\begin{lstlisting}[style=custompython, caption={Finales Testergebnis -- 187 passed, 0 failed}]
$ pytest tests/ -q
.......................................................................
...................................x...x................................
.............................................
187 passed, 2 xfailed, 14 warnings in 36.38s
\end{lstlisting}

\begin{itemize}
      \item \textbf{187 passed}: alle Smoke-, Integrations-, PoC- und Regressionstests
      \item \textbf{2 xfailed}: bekannte Admin-Only-Endpunkte
            (\texttt{test\_list\_organizations\_as\_admin},
            \texttt{test\_create\_organization}), die in der Testumgebung
            ohne Super-Admin-Rechte bewusst nicht verfügbar sind
      \item \textbf{0 failed}: kein Test schlägt fehl
\end{itemize}
